\section{Current Limits and Failure Behavior}\label{sec:limits}

The most important current limit is the event cap in the lowerer.
\texttt{MAX\_EVENTS} is 20,000.
A program that generates 20,001 note events fails during lowering with the diagnostic
\texttt{program exceeds 20000 event limit (has 20001)}.
This is a deliberate safety boundary in the implementation.
It prevents runaway programs from producing very large event streams, but it also means the thesis should not claim that
Motivo currently supports arbitrarily large scores.

The benchmark also revealed a language consistency issue in an old example file.
The current parser does not accept a key-signature header, and key signatures are not part of the current language.
The temporary benchmark copy of the example removed that unsupported line.
The thesis therefore treats key-signature and harmonic-context support as future work, not as implemented behavior.

Single-run memory checks with macOS \texttt{/usr/bin/time -l} reported maximum resident set sizes in the range of
approximately 5.2 MB to 7.8 MB for the largest synthetic cases.
This is useful as an informal local observation, but the time benchmark is the stronger measurement because it was run
with repeated trials.
The memory result should be described cautiously if used outside this chapter.

\subsection{Memory Footprint}\label{detail-subsec:memory-footprint}

The value-oriented design of the AST and IR ensures that the compiler's memory footprint remains minimal.
The use of \texttt{std::variant} supports spatial locality in the CPU cache during AST traversal.
The intermediate representation holds only the necessary atomic data (pitch, start beat, duration, velocity) packed
tightly into \texttt{struct}s. Consequently, even a composition unrolled into large event streams within the
configured implementation limits will consume only a few
megabytes of RAM, allowing the compiler to run seamlessly on low-end hardware or within memory-constrained Docker
containers.
